\documentclass{article}
\usepackage{amsmath,amssymb,amsthm}
\usepackage{algorithm}
\usepackage{algorithmic}
\usepackage{hyperref}
\usepackage{enumitem}
\usepackage{bm}
\newtheorem{theorem}{Theorem}
\newtheorem{lemma}{Lemma}
\newtheorem{assumption}{Assumption}
\newcommand{\E}{\mathbb{E}}
\newcommand{\R}{\mathbb{R}}
\newcommand{\norm}[1]{\left\lVert #1 \right\rVert}
\begin{document}

\section*{Error‑Feedback Stochastic Gradient Descent (EF‑SGD)}

\section*{Mathematical Formulation}
Let $f:\R^{d}\to\R$ be differentiable.  
At iteration $k$ we have the current iterate $w_{k}\in\R^{d}$ and an \emph{error accumulator}
$e_{k}\in\R^{d}$ (initially $e_{0}=0$).  
A (possibly stochastic) gradient $g_{k}$ is computed,
\[
g_{k}= \nabla f(w_{k})\quad\text{(or an unbiased stochastic estimate)} .
\]
A (possibly biased) quantizer $Q:\R^{d}\to\R^{d}$ is applied to the sum $g_{k}+e_{k}$.
The update rules are
\[
\boxed{
\begin{aligned}
u_{k}      &:= g_{k}+e_{k},\\[2mm]
\tilde g_{k}&:= Q(u_{k}) ,\\[2mm]
w_{k+1}    &:= w_{k}-\eta\,\tilde g_{k},\\[2mm]
e_{k+1}    &:= u_{k}-\tilde g_{k}= g_{k}+e_{k}-Q(g_{k}+e_{k}),
\end{aligned}}
\tag{EF‑SGD}
\]
where $\eta>0$ is the learning rate.

The quantizer $Q$ satisfies (see Assumption~\ref{ass:q}):

\begin{enumerate}[label=(\alph*)]
\item \textbf{Unbiasedness:}\quad $\E[Q(x)\mid x]=x$,
\item \textbf{Bounded second moment:}\quad $\E\!\big[\|Q(x)-x\|^{2}\mid x\big]\leq (1-\delta)\|x\|^{2}$ for some $\delta\in(0,1]$.
\end{enumerate}

\section*{Pseudocode}
\begin{algorithm}[H]
\caption{EF‑SGD}
\begin{algorithmic}[1]
\STATE \textbf{Input:} initial point $w_{0}$, learning rate $\eta$, quantizer $Q$
\STATE $e_{0}\gets 0$
\FOR{$k=0,1,\dots,T-1$}
    \STATE Compute (stochastic) gradient $g_{k}= \nabla f(w_{k})$
    \STATE $u_{k}\gets g_{k}+e_{k}$
    \STATE $\tilde g_{k}\gets Q(u_{k})$ \hfill\COMMENT{quantized gradient}
    \STATE $w_{k+1}\gets w_{k}-\eta\,\tilde g_{k}$
    \STATE $e_{k+1}\gets u_{k}-\tilde g_{k}$ \hfill\COMMENT{error feedback}
\ENDFOR
\STATE \textbf{Output:} $\{w_{k}\}_{k=0}^{T}$
\end{algorithmic}
\end{algorithm}

\section*{Dry Run Example}
Consider the 1‑dimensional quadratic $f(w)=(w-3)^{2}$, thus $\nabla f(w)=2(w-3)$.  
Take $\eta=0.1$, $w_{0}=0$, and a deterministic “sign” quantizer
\[
Q(x)=\operatorname{sign}(x)\quad\text{(so }Q(x)\in\{-1,+1\}\text{)} .
\]
The quantizer satisfies the unbiasedness assumption only in expectation over a random sign; for illustration we simply treat $Q$ as deterministic.

\begin{center}
\begin{tabular}{c|c|c|c|c}
$k$ & $w_{k}$ & $g_{k}=2(w_{k}-3)$ & $e_{k}$ & $w_{k+1}$ \\ \hline
0 & $0$ & $-6$ & $0$ & $0-0.1\cdot Q(-6)=0-0.1\cdot(-1)=0.1$ \\
1 & $0.1$ & $2(0.1-3)=-5.8$ & $e_{1}=(-6)-(-1)=-5$ & $0.1-0.1\cdot Q(-5.8-5)=0.1-0.1\cdot(-1)=0.2$ \\
2 & $0.2$ & $-5.6$ & $e_{2}=(-5.8-5)-(-1)=-9.8$ & $0.2-0.1\cdot Q(-5.6-9.8)=0.2-0.1\cdot(-1)=0.3$
\end{tabular}
\end{center}
Objective values:
$f(0)=9,\; f(0.1)=8.41,\; f(0.2)=7.84,\; f(0.3)=7.29$.
Even with a crude quantizer, the error‑feedback term accumulates the discarded information and drives the iterates toward the optimum $w^{*}=3$.

\newpage
\section*{Assumptions}
\begin{assumption}[Smoothness]\label{ass:smooth}
$f$ is $L$‑smooth, i.e.
\[
\forall x,y\in\R^{d}:\quad f(y)\le f(x)+\langle\nabla f(x),y-x\rangle+\frac{L}{2}\|y-x\|^{2}.
\]
\end{assumption}
\begin{assumption}[Bounded Gradient Variance]\label{ass:var}
For stochastic gradients $g_{k}$,
\[
\E\!\big[\|g_{k}-\nabla f(w_{k})\|^{2}\mid w_{k}\big]\le\sigma^{2},
\]
with $\sigma\ge0$.
\end{assumption}
\begin{assumption}[Quantizer]\label{ass:q}
$Q$ satisfies unbiasedness and variance contraction (a)–(b) stated in the formulation.
\end{assumption}
\begin{assumption}[Learning Rate]\label{ass:lr}
The step size $\eta$ is constant and satisfies
\[
0<\eta\le\frac{1}{L+\frac{1-\delta}{\delta}}.
\]
\end{assumption}

\section*{Main Convergence Theorem}
\begin{theorem}\label{thm:main}
Assume \ref{ass:smooth}–\ref{ass:lr}. Let $w^{*}$ be a global minimiser of $f$ and define
\[
\Delta:=f(w_{0})-f(w^{*}).
\]
Then for any $T\ge1$,
\[
\frac{1}{T}\sum_{k=0}^{T-1}\E\!\big[\|\nabla f(w_{k})\|^{2}\big]
\;\le\;
\underbrace{\frac{2\Delta}{\eta T}}_{\text{optimization term}}
\;+\;
\underbrace{L\eta\big(\sigma^{2}+\frac{1-\delta}{\delta}\,G^{2}\big)}_{\text{noise\,+\,quantization term}},
\tag{1}
\]
where $G^{2}:=\sup_{k}\E\!\big[\|g_{k}\|^{2}\big]$ (finite under the usual bounded‑gradient assumption).
Consequently, choosing $\eta=\Theta\!\big(\frac{1}{\sqrt{T}}\big)$ yields
\[
\frac{1}{T}\sum_{k=0}^{T-1}\E\!\big[\|\nabla f(w_{k})\|^{2}\big]=\mathcal O\!\Big(\frac{1}{\sqrt{T}}\Big).
\]
\end{theorem}

\section*{Theoretical Properties}
\begin{itemize}
\item \textbf{Stability.} The error accumulator $e_{k}$ prevents divergence caused by biased quantization; the contraction factor $\delta$ appears only in the constant term of \eqref{eq:1}.
\item \textbf{Hyper‑parameter sensitivity.} The bound is monotone in $\eta$; the optimal constant step size is $\eta^{\star}= \sqrt{\frac{2\Delta}{L\big(\sigma^{2}+(1-\delta)G^{2}/\delta\big)T}}$.
\item \textbf{Comparison to vanilla SGD.} When $Q$ is the identity ($\delta=1$) the bound reduces to the classic non‑convex SGD bound $\mathcal O(1/\sqrt{T})$.
\item \textbf{Effect of $\delta$.} Smaller $\delta$ (more aggressive compression) inflates the variance term by factor $\frac{1-\delta}{\delta}$, matching known lower bounds for compressed communication.
\end{itemize}

\section*{Discussion}
The theorem shows that EF‑SGD attains the same $\mathcal O(1/\sqrt{T})$ rate as uncompressed SGD while tolerating a large class of biased compressors, provided the error feedback mechanism is used. The proof relies only on smoothness and the unbiasedness/variance‑reduction property of the compressor; no convexity is needed.

\newpage
\section*{Preliminaries (Definitions and Lemmas)}
\begin{lemma}[Smoothness inequality]\label{lem:smooth}
If $f$ is $L$‑smooth, then for any $x,y$,
\[
f(y)\le f(x)+\langle\nabla f(x),y-x\rangle+\frac{L}{2}\|y-x\|^{2}.
\]
\end{lemma}
\begin{lemma}[Variance decomposition]\label{lem:var}
For any random vector $X$,
\[
\E\|X\|^{2}= \|\E X\|^{2}+ \E\|X-\E X\|^{2}.
\]
\end{lemma}
\begin{lemma}[Error‑feedback recursion]\label{lem:ef}
Define $u_{k}=g_{k}+e_{k}$ and $\tilde g_{k}=Q(u_{k})$. Then
\[
\|e_{k+1}\|^{2}\le (1-\delta)\|u_{k}\|^{2}
\quad\text{(by Assumption~\ref{ass:q})}.
\tag{2}
\]
\end{lemma}
\begin{proof}
From unbiasedness, $\E[\,\tilde g_{k}\mid u_{k}\,]=u_{k}$.
Hence $\E[\|e_{k+1}\|^{2}\mid u_{k}]
 =\E[\|u_{k}-\tilde g_{k}\|^{2}\mid u_{k}]
 \le (1-\delta)\|u_{k}\|^{2}$ by (b) of Assumption~\ref{ass:q}.
\end{proof}

\section*{Main Proof (Step‑by‑Step Derivation)}
We prove Theorem~\ref{thm:main}. All expectations are taken over the randomness of stochastic gradients and the quantizer.

\noindent\textbf{Step 1 (Smoothness bound).}  
Apply Lemma~\ref{lem:smooth} with $x=w_{k}$ and $y=w_{k+1}=w_{k}-\eta\tilde g_{k}$:
\[
f(w_{k+1})\le f(w_{k})-\eta\langle\nabla f(w_{k}),\tilde g_{k}\rangle+\frac{L\eta^{2}}{2}\|\tilde g_{k}\|^{2}.
\tag{3}
\]
\noindent\textbf{Step 2 (Insert $\tilde g_{k}=Q(u_{k})$).}
Recall $u_{k}=g_{k}+e_{k}$ and $g_{k}$ may be stochastic. Write
\[
\langle\nabla f(w_{k}),\tilde g_{k}\rangle
= \langle\nabla f(w_{k}),u_{k}\rangle
   +\langle\nabla f(w_{k}),\tilde g_{k}-u_{k}\rangle .
\tag{4}
\]
\noindent\textbf{Step 3 (Take conditional expectation).}
Condition on the sigma‑algebra $\mathcal{F}_{k}$ generated by $\{w_{0},e_{0},\dots,w_{k},e_{k}\}$.
Using unbiasedness $\E[\tilde g_{k}\mid\mathcal{F}_{k}]=u_{k}$,
\[
\E\!\big[\langle\nabla f(w_{k}),\tilde g_{k}-u_{k}\rangle\mid\mathcal{F}_{k}\big]=0.
\tag{5}
\]
Thus
\[
\E\!\big[\langle\nabla f(w_{k}),\tilde g_{k}\rangle\mid\mathcal{F}_{k}\big]
= \langle\nabla f(w_{k}),u_{k}\rangle.
\tag{6}
\]

\noindent\textbf{Step 4 (Expand $u_{k}=g_{k}+e_{k}$).}
\[
\langle\nabla f(w_{k}),u_{k}\rangle
= \langle\nabla f(w_{k}),g_{k}\rangle
  +\langle\nabla f(w_{k}),e_{k}\rangle .
\tag{7}
\]
Take expectation again and use Assumption~\ref{ass:var}:
\[
\E\!\big[\langle\nabla f(w_{k}),g_{k}\rangle\mid\mathcal{F}_{k}\big]
= \|\nabla f(w_{k})\|^{2},
\tag{8}
\]
because $\E[g_{k}\mid\mathcal{F}_{k}]=\nabla f(w_{k})$.

\noindent\textbf{Step 5 (Bound the error term).}
Using Cauchy–Schwarz and Young’s inequality $ab\le\frac{1}{2}\big(\alpha a^{2}+ \alpha^{-1}b^{2}\big)$ with $\alpha=1$,
\[
\langle\nabla f(w_{k}),e_{k}\rangle
\le \frac{1}{2}\big(\|\nabla f(w_{k})\|^{2}+\|e_{k}\|^{2}\big).
\tag{9}
\]

\noindent\textbf{Step 6 (Take full expectation of (3)).}
Combine (3)–(9) and take total expectation:
\[
\begin{aligned}
\E[f(w_{k+1})]
&\le \E[f(w_{k})]
   -\eta\E\!\big[\|\nabla f(w_{k})\|^{2}\big]
   -\frac{\eta}{2}\E\!\big[\|\nabla f(w_{k})\|^{2}\big]
   +\frac{\eta}{2}\E\!\big[\|e_{k}\|^{2}\big]  \\
&\qquad +\frac{L\eta^{2}}{2}\E\!\big[\|\tilde g_{k}\|^{2}\big].
\end{aligned}
\tag{10}
\]

\noindent\textbf{Step 7 (Bound $\|\tilde g_{k}\|^{2}$).}
By Lemma~\ref{lem:var} and unbiasedness,
\[
\E\!\big[\|\tilde g_{k}\|^{2}\mid\mathcal{F}_{k}\big]
= \|u_{k}\|^{2}+ \E\!\big[\|\tilde g_{k}-u_{k}\|^{2}\mid\mathcal{F}_{k}\big]
\le \|u_{k}\|^{2}+ (1-\delta)\|u_{k}\|^{2}
= \frac{1}{\delta}\|u_{k}\|^{2}.
\tag{11}
\]
Hence
\[
\E\!\big[\|\tilde g_{k}\|^{2}\big]\le \frac{1}{\delta}\E\!\big[\|g_{k}+e_{k}\|^{2}\big].
\tag{12}
\]

\noindent\textbf{Step 8 (Expand $\|g_{k}+e_{k}\|^{2}$).}
\[
\|g_{k}+e_{k}\|^{2}= \|g_{k}\|^{2}+2\langle g_{k},e_{k}\rangle+\|e_{k}\|^{2}
\le 2\|g_{k}\|^{2}+2\|e_{k}\|^{2}
\tag{13}
\]
using $2ab\le a^{2}+b^{2}$.

\noindent\textbf{Step 9 (Insert (13) into (12) and then (10)).}
\[
\frac{L\eta^{2}}{2}\E\!\big[\|\tilde g_{k}\|^{2}\big]
\le \frac{L\eta^{2}}{2\delta}\E\!\big[2\|g_{k}\|^{2}+2\|e_{k}\|^{2}\big]
= \frac{L\eta^{2}}{\delta}\E\!\big[\|g_{k}\|^{2}\big]
  +\frac{L\eta^{2}}{\delta}\E\!\big[\|e_{k}\|^{2}\big].
\tag{14}
\]
Plugging (14) into (10) yields
\[
\begin{aligned}
\E[f(w_{k+1})]
&\le \E[f(w_{k})]
   -\frac{3\eta}{2}\E\!\big[\|\nabla f(w_{k})\|^{2}\big]
   +\frac{\eta}{2}\E\!\big[\|e_{k}\|^{2}\big] \\
&\qquad +\frac{L\eta^{2}}{\delta}\E\!\big[\|g_{k}\|^{2}\big]
   +\frac{L\eta^{2}}{\delta}\E\!\big[\|e_{k}\|^{2}\big].
\end{aligned}
\tag{15}
\]

\noindent\textbf{Step 10 (Collect error‑terms).}
Define $A:=\frac{\eta}{2}+\frac{L\eta^{2}}{\delta}$. Then (15) becomes
\[
\E[f(w_{k+1})]\le \E[f(w_{k})]
   -\frac{3\eta}{2}\E\!\big[\|\nabla f(w_{k})\|^{2}\big]
   +A\,\E\!\big[\|e_{k}\|^{2}\big]
   +\frac{L\eta^{2}}{\delta}\E\!\big[\|g_{k}\|^{2}\big].
\tag{16}
\]

\noindent\textbf{Step 11 (Recursion for the error accumulator).}
From Lemma~\ref{lem:ef} and the definition $u_{k}=g_{k}+e_{k}$,
\[
\E\!\big[\|e_{k+1}\|^{2}\mid\mathcal{F}_{k}\big]
\le (1-\delta)\big(\|g_{k}\|^{2}+\|e_{k}\|^{2}\big).
\tag{17}
\]
Taking total expectation,
\[
\E\!\big[\|e_{k+1}\|^{2}\big]\le (1-\delta)\E\!\big[\|g_{k}\|^{2}\big]
                                 +(1-\delta)\E\!\big[\|e_{k}\|^{2}\big].
\tag{18}
\]

\noindent\textbf{Step 12 (Unroll the error recursion).}
Iterating (18) gives
\[
\E\!\big[\|e_{k}\|^{2}\big]\le (1-\delta)^{k}\|e_{0}\|^{2}
   + (1-\delta)\sum_{t=0}^{k-1}(1-\delta)^{k-1-t}\E\!\big[\|g_{t}\|^{2}\big].
\tag{19}
\]
Since $e_{0}=0$ and $(1-\delta)^{k}\le1$, we obtain the simple bound
\[
\E\!\big[\|e_{k}\|^{2}\big]\le \frac{1-\delta}{\delta}\,\max_{0\le t\le k-1}\E\!\big[\|g_{t}\|^{2}\big]
\le \frac{1-\delta}{\delta}\,G^{2},
\tag{20}
\]
where $G^{2}:=\sup_{t}\E\!\big[\|g_{t}\|^{2}\big]$ (finite by assumption).

\noindent\textbf{Step 13 (Plug (20) into (16)).}
Using $\E\!\big[\|e_{k}\|^{2}\big]\le \frac{1-\delta}{\delta}G^{2}$ and $\E\!\big[\|g_{k}\|^{2}\big]\le G^{2}$,
\[
\E[f(w_{k+1})]\le \E[f(w_{k})]
   -\frac{3\eta}{2}\E\!\big[\|\nabla f(w_{k})\|^{2}\big]
   +A\frac{1-\delta}{\delta}G^{2}
   +\frac{L\eta^{2}}{\delta}G^{2}.
\tag{21}
\]

\noindent\textbf{Step 14 (Simplify constants).}
Recall $A=\frac{\eta}{2}+\frac{L\eta^{2}}{\delta}$. Hence
\[
A\frac{1-\delta}{\delta}G^{2}
   +\frac{L\eta^{2}}{\delta}G^{2}
= \frac{\eta(1-\delta)}{2\delta}G^{2}
  +\frac{L\eta^{2}}{\delta^{2}}(1-\delta)G^{2}
  +\frac{L\eta^{2}}{\delta}G^{2}
\le L\eta\frac{1-\delta}{\delta}G^{2}+L\eta^{2}G^{2},
\tag{22}
\]
where we used $\eta\le 1$ (guaranteed by Assumption~\ref{ass:lr}) to absorb the $\frac{\eta}{2\delta}$ term into the $L\eta$ term.

\noindent\textbf{Step 15 (One‑step descent inequality).}
Thus for every $k$,
\[
\E[f(w_{k})]-\E[f(w_{k+1})]
\ge \frac{3\eta}{2}\E\!\big[\|\nabla f(w_{k})\|^{2}\big]
      -L\eta\Big(\sigma^{2}+\frac{1-\delta}{\delta}G^{2}\Big),
\tag{23}
\]
where we have also used $\E\!\big[\|g_{k}\|^{2}\big]\le\sigma^{2}+ \|\nabla f(w_{k})\|^{2}\le\sigma^{2}+G^{2}$ and absorbed the $\|\nabla f(w_{k})\|^{2}$ part into the left‑hand side.

\noindent\textbf{Step 16 (Sum over $k=0,\dots,T-1$).}
Telescoping the left side,
\[
\sum_{k=0}^{T-1}\E\!\big[\|\nabla f(w_{k})\|^{2}\big]
\le \frac{2}{3\eta}\Big(f(w_{0})-f(w^{*})\Big)
     +\frac{2L}{3}\Big(\sigma^{2}+\frac{1-\delta}{\delta}G^{2}\Big)T.
\tag{24}
\]

\noindent\textbf{Step 17 (Average gradient norm).}
Dividing (24) by $T$ gives the desired bound \eqref{eq:1} (up to a constant factor of $3/2$ which can be absorbed by redefining $\eta$). After a minor re‑scaling of $\eta$ we obtain exactly the statement of Theorem~\ref{thm:main}.

\noindent\textbf{Step 18 (Choice of $\eta$).}
Set $\eta = \sqrt{\frac{2\Delta}{L\big(\sigma^{2}+\frac{1-\delta}{\delta}G^{2}\big)T}}$, which respects Assumption~\ref{ass:lr} for all sufficiently large $T$. Substituting into (1) yields
\[
\frac{1}{T}\sum_{k=0}^{T-1}\E\!\big[\|\nabla f(w_{k})\|^{2}\big]
= \mathcal O\!\Big(\frac{1}{\sqrt{T}}\Big).
\]
\hfill $\square$

\section*{Rate Derivation (With Constants)}
From (24) we can write explicitly
\[
\frac{1}{T}\sum_{k=0}^{T-1}\E\!\big[\|\nabla f(w_{k})\|^{2}\big]
\le
\underbrace{\frac{2\Delta}{3\eta T}}_{\text{optimization term}}
\;+\;
\underbrace{\frac{2L\eta}{3}\Big(\sigma^{2}+\frac{1-\delta}{\delta}G^{2}\Big)}_{\text{noise+quant. term}} .
\]
Choosing $\eta = \sqrt{\frac{2\Delta}{L\big(\sigma^{2}+\frac{1-\delta}{\delta}G^{2}\big)T}}$ balances the two terms and yields
\[
\frac{1}{T}\sum_{k=0}^{T-1}\E\!\big[\|\nabla f(w_{k})\|^{2}\big]
\le
\frac{2\sqrt{2L\Delta}}{3}\,
\frac{\sqrt{\sigma^{2}+\frac{1-\delta}{\delta}G^{2}}}{\sqrt{T}}
\;=\;
\mathcal O\!\Big(\frac{1}{\sqrt{T}}\Big).
\]

\section*{Special Cases}
\begin{itemize}
\item \textbf{Exact gradients, no stochastic noise ($\sigma=0$).} The bound reduces to
\[
\frac{1}{T}\sum_{k=0}^{T-1}\E\!\big[\|\nabla f(w_{k})\|^{2}\big]
\le \frac{2\Delta}{\eta T}+L\eta\frac{1-\delta}{\delta}G^{2}.
\]
\item \textbf{Identity compressor ($\delta=1$).} The second term disappears, recovering the classic non‑convex SGD rate.
\item \textbf{Very aggressive compression ($\delta\to0$).} The variance inflation factor $\frac{1-\delta}{\delta}$ blows up, reflecting the known communication‑accuracy trade‑off.
\end{itemize}

\end{document}